\section{Orbit orders and characters}
\label{V:sec-order-weight}

We use the left action on sections of a linearized line bundle:
$(g\cdot s)(x)=g\bigl(s(g^{-1}x)\bigr)$.  Thus a character
$\alpha$ means $g\cdot s=\alpha(g)s$.  Fixing this convention is
necessary when an order estimate is converted into a weight estimate.
The next two elementary lemmas will be used both for approximate
weights and for exact attainment of a supporting weight.

\begin{lemma}[The order on a generic orbit]\label{V:order-weight}
Let an algebraic torus $T$ act on a smooth integral projective variety
$X$, and let $L$ be a $T$-linearized line bundle.  Let
$\lambda:\mathbb G_m\to T$ be a cocharacter, and let $S$ be its
source component: a dense open subset of $X$ tends to $S$ as the
parameter tends to zero.  Let $\mu$ be the $\lambda$-weight on
$L|_S$.  If $0\ne s\in H^0(X,L)$ has $\lambda$-weight $w$, then
\[
 v_\lambda(\operatorname{div}s)=\mu-w,
\]
where $v_\lambda$ is the order on the arc
$t\mapsto\lambda(t)\xi$ and $\xi$ is the generic point, viewed as
a $\mathbb C(X)$-valued point.  For the zero cocharacter the order
and both weights are zero.
\end{lemma}

\begin{proof}
The orbit map is defined over $\mathbb C(X)(t)$ and remains generic
on $X$.  By properness it extends to the formal disc over
$\mathbb C(X)$.  Its center belongs to $S$.  Pull back $L$ to this disc.
If the pulled-back section has order $d$, its first nonzero class
belongs to the fiber of $L$ tensored with the degree-$d$ parameter
ideal quotient.  Parameter scaling acts on this quotient, under
the declared action on sections, with weight $-d$.  The first
nonzero class therefore has weight $\mu-d$.  Equivariance gives
$w=\mu-d$.  The zero-cocharacter arc is constant at the generic
point; every nonzero rational function is a unit there, which
proves the last assertion directly.
\end{proof}

\begin{lemma}[Multiplying characters]\label{V:convex-product}
Let $T$ act on an integral projective variety $X$, and let $L$ be
a $T$-linearized line bundle.  Suppose, for $1\le i\le N$, that $s_i\in H^0(X,r_iL)$ are
nonzero eigensections with characters $\alpha_i$, where
$r_i\in\mathbb Z_{>0}$.
If zero lies in the convex hull of the rational points
$\alpha_i/r_i$, then $H^0(X,rL)^T\ne0$ for some $r>0$.
\end{lemma}

\begin{proof}
The defining linear equations and inequalities have rational
coefficients, so there are rational numbers $c_i\ge0$ satisfying
$\sum_i c_i=1$ and $\sum_i c_i\alpha_i/r_i=0$.  Clear the
denominators of $c_i/r_i$ to obtain integers $n_i\ge0$, not all
zero, with $\sum_i n_i\alpha_i=0$.  The product
$\prod_i s_i^{n_i}$ has weight zero and belongs to
$H^0(X,rL)$, where $r=\sum_i n_i r_i>0$.  Its factors are
nonzero at the generic point of the integral variety, so the
product is nonzero.
\end{proof}
